\subsection{Types of quadratic forms.}

Given a quadratic form Q (§ 3, n° 4) on a vector space E over A, we shall say that E is the space of definition of Q and that dim(E) is the dimension of Q. Given two quadratic forms Q, Q' on vector spaces E, E' over A, we shall denote by Q τ Q' their direct sum (§ 3, n° 4). Recall that the direct sum of two neutral forms is neutral (§ 4, n° 2).

Introduce the following relation:

“Q and Q' are non-degenerate quadratic forms of finite dimensions over A, and there exist neutral quadratic forms, N, N' such that Q τ N is equivalent to Q' τ N'”.

This relation, which we shall denote by \(Q \sim Q'\) , is manifestly reflexive and symmetric. It is also transitive: indeed, if \(Q\) , \(Q'\) , \(Q''\) are quadratic forms such that \(Q \sim Q'\) and \(Q' \sim Q''\) , there exist neutral quadratic forms \(M\) , \(M'\) , \(N\) , \(N'\) such that \(Q \tau M\) is equivalent to \(Q' \tau M'\) and \(Q' \tau N\) is equivalent to \(Q'' \tau N'\) ; then \(Q \tau (M \tau N)\) is equivalent to \((Q \tau M) \tau N\) , hence to \((Q' \tau M') \tau N\) , and also to \((Q' \tau N) \tau M'\) , hence again to \((Q'' \tau N') \tau M'\) and to \(Q'' \tau (N' \tau M')\) ; since \(M \tau N\) and \(N' \tau M'\) are neutral, one indeed has \(Q \sim Q''\) . The relation \(Q \sim Q'\) is therefore an equivalence relation between \(Q\) and \(Q'\) . It is clear that, if \(Q\) and \(Q'\) are two non-degenerate quadratic forms of finite dimensions and equivalent, one has \(Q \sim Q'\) .

For every quadratic form Q on A, non-degenerate and of finite dimension, we set

\[
\theta (Q) = \tau_ {x} (X \sim Q),\tag{1}
\]

and we shall say that \(\theta(Q)\) is the type of \(Q\) . If \(Q\) and \(Q'\) are two quadratic forms on \(A\) , non-degenerate and of finite dimensions, the relations \(Q \sim Q'\) and \(\theta(Q) = \theta(Q')\) are equivalent.

PROPOSITION 1. --- Let Q and Q' be two nondegenerate quadratic forms on A, of finite dimension. For Q and Q' to be equivalent, it is necessary and sufficient that they have the same dimension and the same type.

The condition is obviously necessary. Suppose it is satisfied. Then there exist neutral forms N, N′ such that Q ⊥ N and Q′ ⊥ N′ are equivalent. Since these two forms have the same dimension, so do N and N′, which are consequently equivalent (§ 4, no. 2, cor. 2 of prop. 2). Hence Q and Q′ are equivalent by virtue of Witt's theorem (§ 4, no. 3, cor. 1 of th. 1).

PROPOSITION 2. --- The relation “there exists a nondegenerate quadratic form Q on A, of finite dimension, such that X = θ(Q)” is collectivizing in X (Set Theory, Chap. II, § 1, no. 4).

Let indeed V be a vector space of infinite dimension over A, S the set of nondegenerate quadratic forms defined on the finite-dimensional subspaces of V, and W the set of \(\theta(Q)\) for \(Q \in \mathfrak{S}\) . It is clear that every nondegenerate quadratic form \(Q'\) of finite dimension over A is equivalent to at least one element of \(\mathfrak{S}\) ; whence \(\theta(Q') \in \mathfrak{W}\) , which proves our assertion.

